\documentclass[a4paper]{article}
\usepackage{ctex}
\ctexset{proofname=\heiti{证明}}
\usepackage{amsmath,amssymb,amsthm}
\usepackage{mathtools}
\usepackage{bm}
\usepackage{enumitem}
\usepackage{booktabs}
\usepackage{array}
\usepackage{graphicx}

\IfFileExists{iidef.sty}{
    \usepackage[thehwcnt = 4]{iidef}
    \thecourseinstitute{清华大学}
    \thecoursename{人工智能原理}
    \theterm{2026年春季学期}
    \hwname{作业}
    \slname{\heiti{解}}
}{
    \makeatletter
    \newcommand{\thecourseinstitute}[1]{\def\@courseinstitute{#1}}
    \newcommand{\thecoursename}[1]{\def\@coursename{#1}}
    \newcommand{\theterm}[1]{\def\@term{#1}}
    \newcommand{\hwname}[1]{\def\@hwname{#1}}
    \newcommand{\slname}[1]{\def\@slname{#1}}
    \thecourseinstitute{清华大学}
    \thecoursename{人工智能原理}
    \theterm{2026年春季学期}
    \hwname{作业}
    \slname{\heiti{解}}
    \newcommand{\courseheader}{
        \begin{center}
            {\Large \heiti \@courseinstitute \quad \@coursename}\\[0.5em]
            {\large \@term \quad \@hwname 4 \quad \@slname}
        \end{center}
        \vspace{1em}
    }
    \newcommand{\name}[1]{\noindent #1 \par\vspace{1em}}
    \newenvironment{solution}{\par\noindent\textbf{解：}\par}{\par}
    \makeatother
}

\begin{document}
\courseheader
\name{姓名：杜一郎\qquad 学号：2023011618}

\begin{enumerate}
\setlength{\itemsep}{3\parskip}

\item 线性回归

\begin{solution}
数据为
\[
x=[6,8,10,12,13,14,15,16],
\]
\[
y=[68,73,82,85,87,91,92,98].
\]
使用最小二乘拟合
\[
\hat y=wx+b.
\]
设样本均值为 \(\bar x,\bar y\)，则
\[
w=\frac{\sum_i (x_i-\bar x)(y_i-\bar y)}{\sum_i (x_i-\bar x)^2},
\qquad
b=\bar y-w\bar x.
\]
代入数据得到
\[
w=2.8304,\qquad b=51.2427.
\]
所以拟合直线为
\[
\hat y=2.8304x+51.2427.
\]

预测值如下：
\[
\begin{array}{c|c|c}
\toprule
x & y & \hat y\\
\midrule
6 & 68 & 68.225\\
8 & 73 & 73.886\\
10 & 82 & 79.547\\
12 & 85 & 85.208\\
13 & 87 & 88.038\\
14 & 91 & 90.868\\
15 & 92 & 93.699\\
16 & 98 & 96.529\\
\bottomrule
\end{array}
\]

平均绝对误差为
\[
\mathrm{MAE}
=\frac1n\sum_i |y_i-\hat y_i|
=1.0139.
\]
均方误差为
\[
\mathrm{MSE}
=\frac1n\sum_i (y_i-\hat y_i)^2
=1.6301.
\]
\end{solution}

\item Logistic 回归

\begin{solution}
Logistic 回归模型为
\[
z=w^\top x+b,
\]
\[
h_{w,b}(x)=\sigma(z)=\frac{1}{1+e^{-z}}.
\]
模型输出 \(\hat y=h_{w,b}(x)\)。若需要二分类标签，可取阈值 \(0.5\)：当 \(\hat y\ge 0.5\) 时判为 1，否则判为 0。

单样本交叉熵损失为
\[
L(w,b)
=-\left[y\log h+(1-y)\log(1-h)\right],
\]
其中 \(h=h_{w,b}(x)\)。又有
\[
\sigma'(z)=\sigma(z)(1-\sigma(z))=h(1-h).
\]
先对 \(z\) 求导：
\[
\frac{\partial L}{\partial z}
=
\frac{\partial L}{\partial h}\frac{\partial h}{\partial z}.
\]
其中
\[
\frac{\partial L}{\partial h}
=-\frac{y}{h}+\frac{1-y}{1-h},
\]
所以
\[
\frac{\partial L}{\partial z}
=
\left(-\frac{y}{h}+\frac{1-y}{1-h}\right)h(1-h)
=h-y.
\]
由于
\[
\frac{\partial z}{\partial w_i}=x_i,\qquad
\frac{\partial z}{\partial b}=1,
\]
得到
\[
\frac{\partial L}{\partial w_i}=(h-y)x_i,
\]
\[
\frac{\partial L}{\partial b}=h-y.
\]
用学习率 \(\alpha>0\) 做梯度下降，更新公式为
\[
w_i\leftarrow w_i-\alpha(h-y)x_i,
\]
\[
b\leftarrow b-\alpha(h-y).
\]
\end{solution}

\end{enumerate}
\end{document}
